% GENERATED FILE. Edit canonical lessons, then run npm run generate:books.
% canonical-source-hash: fnv1a64:88d8b64f11b6ce97
% canonical-lessons: IT-C408-benzinaio, IT-C408-gasolio, IT-C408-carburante, IT-C408-serbatoio, IT-C408-rifornire

\chapter{Dal benzinaio}
\label{ch:it-dal-benzinaio}
\hlchaptermodality{408}

\begin{chapteropening}
\textbf{By the end of this chapter:} \emph{I can fill up the car: the attendant, diesel, fuel, the tank and refuelling.}

Complete the last lesson of chapter 408: I can fill up the car: the attendant, diesel, fuel, the tank and refuelling.
\end{chapteropening}

\section[il benzinaio]{il benzinaio — the petrol station attendant; the petrol station}

\label{lesson:IT-C408-benzinaio}



\begin{quote}
Before the new one: say the Italian for the horn, then the Italian for to mist up.
\end{quote}

\subsection*{You'll want to know: il benzinaio}
\textbf{il benzinaio} — \textquotedblleft{}the petrol station attendant; the petrol station\textquotedblright{}.

\begin{quote}
Il benzinaio controlla l'olio.
\end{quote}

\begin{quote}
\emph{The attendant checks the oil.}
\end{quote}

In speech \emph{andare dal benzinaio} means to go to the petrol station.

\begin{cousinweb}
\emph{benzinaio} is built on \emph{benzina} 'petrol', which is usually traced to German \emph{Benzin}, named after benzoin, a resin; the word came into Italian as a borrowing.
\end{cousinweb}

\begin{grammarlens}[title={What you've built}]
One more word for telling, explaining and describing at B1.
\end{grammarlens}

\subsection*{Your turn}
\begin{itemize}
\raggedright
  \item \emph{Say it:} \emph{il benzinaio}
  \item \emph{Say it:} \emph{il benzinaio}, once more
  \item \emph{Say it:} a sentence of your own with \emph{il benzinaio}
  \item \emph{Recall:} say the Italian for the horn, then the Italian for to mist up, then say \emph{il benzinaio} again
\end{itemize}

\begin{tcolorbox}[breakable,colback=teal!4,colframe=teal!35!black,title={Before you move on}]
What does il benzinaio mean? (\textquotedblleft{}The petrol station attendant; the petrol station\textquotedblright{}.) Say it once more.
\end{tcolorbox}
\section[il gasolio]{il gasolio — the diesel (fuel)}

\label{lesson:IT-C408-gasolio}



\begin{quote}
Before the new one: say the Italian for to mist up, then the Italian for the petrol station attendant.
\end{quote}

\subsection*{You'll want to know: il gasolio}
\textbf{il gasolio} — \textquotedblleft{}the diesel (fuel)\textquotedblright{}.

\begin{quote}
La mia macchina va a gasolio.
\end{quote}

\begin{quote}
\emph{My car runs on diesel.}
\end{quote}

Also \emph{il diesel}. \emph{Andare a gasolio} — to run on diesel.

\begin{cousinweb}
\emph{gasolio} is built from \emph{gas}, a word coined in Dutch from Greek \emph{kháos} 'empty space', and \emph{olio} 'oil', which is inherited from Latin \emph{oleum} (EtymDB).
\end{cousinweb}

\begin{grammarlens}[title={What you've built}]
One more word for telling, explaining and describing at B1.
\end{grammarlens}

\subsection*{Your turn}
\begin{itemize}
\raggedright
  \item \emph{Say it:} \emph{il gasolio}
  \item \emph{Say it:} \emph{il gasolio}, once more
  \item \emph{Say it:} a sentence of your own with \emph{il gasolio}
  \item \emph{Recall:} say the Italian for to mist up, then the Italian for the petrol station attendant, then say \emph{il gasolio} again
\end{itemize}

\begin{tcolorbox}[breakable,colback=teal!4,colframe=teal!35!black,title={Before you move on}]
What does il gasolio mean? (\textquotedblleft{}The diesel (fuel)\textquotedblright{}.) Say it once more.
\end{tcolorbox}
\section[il carburante]{il carburante — the fuel}

\label{lesson:IT-C408-carburante}



\begin{quote}
Before the new one: say the Italian for the petrol station attendant, then the Italian for the diesel.
\end{quote}

\subsection*{You'll want to know: il carburante}
\textbf{il carburante} — \textquotedblleft{}the fuel\textquotedblright{}.

\begin{quote}
Il carburante costa sempre di più.
\end{quote}

\begin{quote}
\emph{Fuel costs more and more.}
\end{quote}

The general word: \emph{la benzina} and \emph{il gasolio} are both \emph{carburante}.

\begin{cousinweb}
\emph{carburante} comes from \emph{carburare} 'to mix fuel with air', which was borrowed from French \emph{carburer}, built on Latin \emph{carbō} 'charcoal' (Lewis \& Short).
\end{cousinweb}

\begin{grammarlens}[title={What you've built}]
One more word for telling, explaining and describing at B1.
\end{grammarlens}

\subsection*{Your turn}
\begin{itemize}
\raggedright
  \item \emph{Say it:} \emph{il carburante}
  \item \emph{Say it:} \emph{il carburante}, once more
  \item \emph{Say it:} a sentence of your own with \emph{il carburante}
  \item \emph{Recall:} say the Italian for the petrol station attendant, then the Italian for the diesel, then say \emph{il carburante} again
\end{itemize}

\begin{tcolorbox}[breakable,colback=teal!4,colframe=teal!35!black,title={Before you move on}]
What does il carburante mean? (\textquotedblleft{}The fuel\textquotedblright{}.) Say it once more.
\end{tcolorbox}
\section[il serbatoio]{il serbatoio — the tank (for fuel); the reservoir}

\label{lesson:IT-C408-serbatoio}



\begin{quote}
Before the new one: say the Italian for the diesel, then the Italian for the fuel.
\end{quote}

\subsection*{You'll want to know: il serbatoio}
\textbf{il serbatoio} — \textquotedblleft{}the tank (for fuel); the reservoir\textquotedblright{}.

\begin{quote}
Il serbatoio è quasi vuoto.
\end{quote}

\begin{quote}
\emph{The tank is almost empty.}
\end{quote}

\emph{Fare il pieno} — to fill the tank.

\begin{cousinweb}
\emph{serbatoio} is built on \emph{serbare} 'to keep, to store', which is inherited from Latin \emph{servāre} 'to keep safe' (EtymDB).
\end{cousinweb}

\begin{grammarlens}[title={What you've built}]
One more word for telling, explaining and describing at B1.
\end{grammarlens}

\subsection*{Your turn}
\begin{itemize}
\raggedright
  \item \emph{Say it:} \emph{il serbatoio}
  \item \emph{Say it:} \emph{il serbatoio}, once more
  \item \emph{Say it:} a sentence of your own with \emph{il serbatoio}
  \item \emph{Recall:} say the Italian for the diesel, then the Italian for the fuel, then say \emph{il serbatoio} again
\end{itemize}

\begin{tcolorbox}[breakable,colback=teal!4,colframe=teal!35!black,title={Before you move on}]
What does il serbatoio mean? (\textquotedblleft{}The tank (for fuel); the reservoir\textquotedblright{}.) Say it once more.
\end{tcolorbox}
\section[rifornire]{rifornire — to refuel; to supply}

\label{lesson:IT-C408-rifornire}



\begin{quote}
Before the new one: say the Italian for the fuel, then the Italian for the tank.
\end{quote}

\subsection*{You'll want to know: rifornire}
\textbf{rifornire} — \textquotedblleft{}to refuel; to supply\textquotedblright{}.

\begin{quote}
Dobbiamo rifornire la macchina prima del viaggio.
\end{quote}

\begin{quote}
\emph{We need to refuel the car before the trip.}
\end{quote}

\emph{Il rifornimento} is refuelling. Conjugated like \emph{finire}: \emph{rifornisco}.

\begin{cousinweb}
\emph{rifornire} is built from \emph{ri-} 'again' and \emph{fornire} 'to supply', which is usually traced to a Germanic verb; it was borrowed into Italian through Old French \emph{fornir}.
\end{cousinweb}

\begin{grammarlens}[title={What you've built}]
One more word for telling, explaining and describing at B1.
\end{grammarlens}

\subsection*{Your turn}
\begin{itemize}
\raggedright
  \item \emph{Say it:} \emph{rifornire}
  \item \emph{Say it:} \emph{rifornire}, once more
  \item \emph{Say it:} a sentence of your own with \emph{rifornire}
  \item \emph{Recall:} say the Italian for the fuel, then the Italian for the tank, then say \emph{rifornire} again
\end{itemize}

\begin{tcolorbox}[breakable,colback=teal!4,colframe=teal!35!black,title={Before you move on}]
What does rifornire mean? (\textquotedblleft{}To refuel; to supply\textquotedblright{}.) Say it once more.
\end{tcolorbox}
